\section{Calibration and Empirical Validation}\label{sec:calibration}

\subsection{Incident data}
We compiled the packages compromised in both Shai-Hulud waves from public indicator lists and vendor
reports~\citep{stepsecurity2025shaihulud,socket2025shaihulud,wiz2025shaihulud2,datadog2025shaihulud2}.
This gives \ValOneListed{} packages for wave~1 and \ValTwoListed{} for wave~2. Publish timestamps come
from registry metadata, which survives takedown. Three facts from the incidents anchor the calibration:
\begin{enumerate}
  \item \textbf{Detection lag.} Wave~1 went about 24\,h from the first malicious publish to the first
        public signal. Wave~2 was flagged by automated feeds within minutes.
  \item \textbf{Plateau.} 90\% of the compromised versions had appeared 37\,h (wave~1) and 13.4\,h
        (wave~2) after the first publish.
  \item \textbf{Mechanism.} Most compromised packages were produced by stolen-token republishing.
        Publish bursts match the per-token caps, and only about 230 of the about 800 wave-2 packages have
        any dependents at all. We therefore target a credential-channel share of 75\%.
\end{enumerate}

\subsection{Calibration}
We grid-search $\beta\in[0.005,0.3]$, $\beta_m\in[0.005,0.5]$ and $\gamma\in[0.05,0.4]$ on the main graph,
with each wave's cap and detection tick fixed. Each configuration gets 40 runs, each from a uniformly
random patient zero. Plateau time and channel share are evaluated on major outbreaks only, since the
observed waves are by definition outbreaks that took off. The loss is the squared log-error of the
median $T_{90}$ plus the squared, scaled error of the channel share.

Both best fits lie inside the grid:
\begin{itemize}
  \item \textbf{SH1-calibrated:} $\beta=\CalOneBeta$, $\beta_m=\CalOneBetaM$, $\gamma=\CalOneGamma$,
        giving a plateau at $T_{90}=\CalOneTninety$\,h with \CalOneShare{} of infections via the
        credential channel.
  \item \textbf{SH2-calibrated:} $\beta=\CalTwoBeta$, $\beta_m=\CalTwoBetaM$, $\gamma=\CalTwoGamma$,
        giving $T_{90}=\CalTwoTninety$\,h and a credential share of \CalTwoShare.
\end{itemize}

The data constrain transmission well. Configurations within a loss tolerance of 0.1 of the optimum span
$\beta\in\CalOneBetaRange$ and $\beta_m\in\CalOneBetaMRange$ for wave~1, and $\beta\in\CalTwoBetaRange$
and $\beta_m\in\CalTwoBetaMRange$ for wave~2. They do not identify $\gamma$, since every value we tried
fits, nor the absolute outbreak size, since median major-outbreak sizes range over \CalOneSizeRange{}
packages for wave~1.

We therefore treat outbreak sizes in the high-impact core as \emph{projections}. They answer the
question ``what if a worm with Shai-Hulud's dynamics reached the core?'', not ``what did Shai-Hulud
do''.

\subsection{Validation of the channel structure}
If the credential channel dominates, victims should cluster by shared maintainer much more than by
dependency links. We fetched current registry metadata, never package contents, for
\ValOneN{} of the wave-1 and \ValTwoN{} of the wave-2 packages. We then compared them with random sets of
equal size drawn from the full closure (\cref{tab:validation}).

\begin{table}[t]
\centering
\caption{Clustering of real Shai-Hulud victims. A package counts as \emph{co-maintained} if it shares a
maintainer with another victim, and as \emph{dependency-linked} if it has an installed dependency edge to
another victim. The null is the mean over 200 random sets of equal size from the full closure. The last
column counts victims present in the full-closure graph.}
\label{tab:validation}
\begin{tabular}{lrrrrrr}
\toprule
 & & \multicolumn{2}{c}{Co-maintained} & \multicolumn{2}{c}{Dependency-linked} & In core \\
\cmidrule(lr){3-4}\cmidrule(lr){5-6}
Wave & $n$ & observed & null & observed & null & graph \\
\midrule
Shai-Hulud 1 & 189 & 98\% & 49\% & 32\% & 5\% & 2 \\
Shai-Hulud 2 & 732 & 95\% & 69\% & 23\% & 14\% & 7 \\
\bottomrule
\end{tabular}

\end{table}

The pattern is clear:
\begin{itemize}
  \item \ValOneCoMaint{} of wave-1 victims share a maintainer with another victim, against
        \ValOneCoMaintNull{} in the null. For wave~2 the figures are \ValTwoCoMaint{} and
        \ValTwoCoMaintNull.
  \item Installed dependency links are much rarer: \ValOneDepLinked{} and \ValTwoDepLinked{} of victims.
  \item Almost no victims belong to the high-impact core: \ValOneInFull{} wave-1 and \ValTwoInFull{}
        wave-2 packages appear in the full closure.
\end{itemize}

The null is drawn from the high-impact closure, where monorepo organisations inflate co-maintenance, so
the contrast is conservative. Taken together, the evidence shows that the real waves spread through
stolen tokens among the periphery's small maintainers. Our experiments ask what happens when such a worm
enters the core.
